\documentclass[12pt]{article}
\usepackage[utf8]{inputenc}
\usepackage[T1]{fontenc}
\usepackage{hyperref}
\title{Welaun Language Resources for K-12 Curriculum Development}
\author{Compiled by Claude Sonnet 5.5 (claude-sonnet-5-5) for Andr\'as Kornai}
\date{}
\begin{document}
\maketitle

\section{Introduction}
Welaun (also Bekais, Wekais; ISO 639-3 wlh) is an Austronesian language of the Central Timor group, related to Kemak, Tokodede and Mambae, spoken in Timor-Leste and Indonesian West Timor. Estimates of speakers range from about 2,000 to 5,575 (Edwards). Joshua Project describes it as a first language of the whole ethnic community, and it is not known to be taught in schools.

\section{Primary Educational Resources}
No primers or textbooks were found in search results.

\section{Linguistic and Reference Materials}
\begin{itemize}
\item Owen Edwards, ``Reintroducing Welaun,'' \emph{Oceanic Linguistics} (2019), with a phonology and morphology sketch, a glossed text and two wordlists. The Leiden repository link found in search results is probably this item: \url{https://scholarlypublications.universiteitleiden.nl/access/item%3A2979444/view}
\item Wikipedia, Welaun language: \url{https://en.wikipedia.org/wiki/Welaun_language}
\item ASJP wordlist entry for Welaun: \url{https://asjp.clld.org/languages/WELAUN}
\item Guilherme ``Puru-Berliku'' da Silva, Welaun--Tetun dictionary (2012, about 189 pages), reported as not widely available outside Timor. No link was found.
\item Geoffrey Hull (2003), a description written in Tetun Dili. No link was found.
\end{itemize}

\section{Religious and Community Resources}
\begin{itemize}
\item Joshua Project profile, Welaun in Timor-Leste: \url{https://joshuaproject.net/index.php/people_groups/21838/TT}
\end{itemize}

\section{Digital and Audio Resources}
\begin{itemize}
\item PARADISEC: 2017 recording of the narrative ``Life in Oele'u'' by Antonius Mali, recorded by Owen Edwards, with transcript PDF and WAV and MP3 files: \url{https://catalog.paradisec.org.au/repository/OE8/Welaun20171204_OeleuStory}
\end{itemize}

\section{Curriculum Development Context}
Welaun is spoken in an area where Tetun and Indonesian dominate in schooling. Materials are thin: one modern descriptive article, a Welaun--Tetun dictionary, and a small audio archive. A curriculum would need new orthography work and materials development with the community, using Tetun as the likely bridge language.
\end{document}